% Options for packages loaded elsewhere
\PassOptionsToPackage{unicode}{hyperref}
\PassOptionsToPackage{hyphens}{url}
\PassOptionsToPackage{dvipsnames,svgnames,x11names}{xcolor}
\documentclass[
  10pt,
  fontset=fandol]{ctexart}
\usepackage{xcolor}
\usepackage[margin=16mm]{geometry}
\usepackage{amsmath,amssymb}
\setcounter{secnumdepth}{-\maxdimen} % remove section numbering
\usepackage{iftex}
\ifPDFTeX
  \usepackage[T1]{fontenc}
  \usepackage[utf8]{inputenc}
  \usepackage{textcomp} % provide euro and other symbols
\else % if luatex or xetex
  \usepackage{unicode-math} % this also loads fontspec
  \defaultfontfeatures{Scale=MatchLowercase}
  \defaultfontfeatures[\rmfamily]{Ligatures=TeX,Scale=1}
\fi
\usepackage{lmodern}
\ifPDFTeX\else
  % xetex/luatex font selection
\fi
% Use upquote if available, for straight quotes in verbatim environments
\IfFileExists{upquote.sty}{\usepackage{upquote}}{}
\IfFileExists{microtype.sty}{% use microtype if available
  \usepackage[]{microtype}
  \UseMicrotypeSet[protrusion]{basicmath} % disable protrusion for tt fonts
}{}
\makeatletter
\@ifundefined{KOMAClassName}{% if non-KOMA class
  \IfFileExists{parskip.sty}{%
    \usepackage{parskip}
  }{% else
    \setlength{\parindent}{0pt}
    \setlength{\parskip}{6pt plus 2pt minus 1pt}}
}{% if KOMA class
  \KOMAoptions{parskip=half}}
\makeatother
\setlength{\emergencystretch}{3em} % prevent overfull lines
\providecommand{\tightlist}{%
  \setlength{\itemsep}{0pt}\setlength{\parskip}{0pt}}
\usepackage{amsmath,amssymb,mathtools,needspace}
\xeCJKDeclareCharClass{CJK}{"2032,"0400->"04FF,"2160->"216B,"2460->"2473,"25A0,"25C7,"25CB,"25CF}
\IfFontExistsTF{Arial Unicode MS}{\xeCJKsetup{AutoFallBack=true}\setCJKfallbackfamilyfont{\CJKrmdefault}{Arial Unicode MS}}{}
\setcounter{secnumdepth}{0}
\setlength{\emergencystretch}{2em}
\allowdisplaybreaks
\usepackage{bookmark}
\IfFileExists{xurl.sty}{\usepackage{xurl}}{} % add URL line breaks if available
\urlstyle{same}
\hypersetup{
  pdftitle={内积空间},
  colorlinks=true,
  linkcolor={Maroon},
  filecolor={Maroon},
  citecolor={Blue},
  urlcolor={Blue},
  pdfcreator={LaTeX via pandoc}}

\title{内积空间}
\author{}
\date{}

\begin{document}
\maketitle

\subsubsection{3.3.1 内积空间}\label{ux5185ux79efux7a7aux95f4}

定义3.4 设在区间 \((a,b)\) 内, 非负函数 \(\rho(x)\) 满足以下条件, 就称 \(\rho(x)\) 为区间 \((a,b)\) 内的权函数:

\(1^\circ \int_a^b |x|^n \rho(x) \mathrm{d}x (n = 0,1,\cdots)\) 存在;

\(2^\circ\) 对于非负的连续函数 \(g(x)\), 若 \[
\int_a^b g(x) \rho(x) \mathrm{d}x = 0, \tag{3.3.2}
\] 则在 \((a,b)\) 内 \(g(x) \equiv 0\).

定义3.5 设 \(f(x), g(x) \in C[a,b], \rho(x)\) 是 \([a,b]\) 上的权函数, 积分 \[
(f,g) = \int_a^b \rho(x) f(x) g(x) \mathrm{d}x \tag{3.3.3}
\] 称为函数 \(f(x)\) 与 \(g(x)\) 在 \([a,b]\) 上的内积.

容易验证这样定义的内积满足下列四条公理:

\begin{enumerate}
\def\labelenumi{(\arabic{enumi})}
\item
  \((f,g) = (g,f)\);
\item
  \((cf,g) = c(f,g), c\) 为常数;
\item
  \((f_1 + f_2, g) = (f_1, g) + (f_2, g)\);
\item
  \((f,f) \geqslant 0\), 当且仅当 \(f=0\) 时 \((f,f)=0\).
\end{enumerate}

满足内积定义的函数空间称为内积空间, 因此, 连续函数空间 \(C[a,b]\) 上定义了内积就形成一个内积空间. 这里定义的内积是 \(n\) 维 Euclid 空间 \(\mathbf{R}^n\) 中两个向量内积定义的推广, 设 \[
\boldsymbol{f} = (f_1, f_2, \cdots, f_n)^{\mathrm{T}}, \quad \boldsymbol{g} = (g_1, g_2, \cdots, g_n)^{\mathrm{T}},
\]

\href{../../source-images/pdf-066.jpeg}{查看原页}

其内积定义为 \((\boldsymbol{f}, \boldsymbol{g}) = \sum_{k=1}^{n} f_{k} g_{k}\); 向量 \(\boldsymbol{f} \in \mathbf{R}^{n}\) 的模(范数)定义为 \[
\| \boldsymbol{f} \|_{2} = \left( \sum_{k=1}^{n} f_{k}^{2} \right)^{\frac{1}{2}},
\] 将它推广到任何内积空间中就有下面的定义.

定义 3.6 \(f(x) \in C[a, b]\), 量 \[
\| f \|_{2} = \sqrt{\int_{a}^{b} \rho(x) f^{2}(x) \mathrm{d}x} = \sqrt{(f, f)} \tag{3.3.4}
\]

称为 \(f(x)\) 的 Euclid 范数.

它同样满足范数三条性质(见 3.1 节), 头两条是显然的, 下面定理包含了第三条性质和其他重要结论.

定理 3.5 对于任何 \(f, g \in C[a, b]\), 下列结论成立: \(1^{\circ}\)

\[
|(f,g)|\leqslant\|f\|_2\|g\|_2\qquad\text{(Cauchy-Schwarz 不等式)};\tag{3.3.5}
\] \(2^{\circ}\)

\[
\|f+g\|_2\leqslant\|f\|_2+\|g\|_2\qquad\text{(三角不等式)};\tag{3.3.6}
\] \(3^{\circ} \| f + g \|_{2}^{2} + \| f - g \|_{2}^{2} = 2(\| f \|_{2}^{2} + \| g \|_{2}^{2})\) (平行四边形定律).

证明 若 \(g=0\), 则式(3.3.5)显然成立, 现考虑 \(g \neq 0\). 对于任何实数 \(\lambda\), 有 \[
0 \leqslant (f + \lambda g, f + \lambda g) = (f, f) + 2\lambda(f, g) + \lambda^{2}(g, g).
\]

现取 \(\lambda = -\frac{(f, g)}{\| g \|_{2}^{2}}\), 代入上式得 \[
\| f \|_{2}^{2} - 2 \frac{|(f, g)|^{2}}{\| g \|_{2}^{2}} + \frac{|(f, g)|^{2}}{\| g \|_{2}^{2}} \geqslant 0, \quad |(f, g)|^{2} \leqslant \| f \|_{2}^{2} \| g \|_{2}^{2}.
\] 两边开平方即得式(3.3.5).

利用式(3.3.5)考虑 \[
\| f + g \|_{2}^{2} = (f + g, f + g) = (f, f) + 2(f, g) + (g, g)
\] \[
\leqslant \| f \|_{2}^{2} + 2 |(f, g)| + \| g \|_{2}^{2}
\] \[
\leqslant \| f \|_{2}^{2} + 2 \| f \|_{2} \| g \|_{2} + \| g \|_{2}^{2}
\] \[
= (\| f \|_{2} + \| g \|_{2})^{2},
\] 两边开方则得式(3.3.6).

平行四边形定律可直接计算得出, 可作为习题由读者完成. 证毕.

在 \(n\) 维空间中两个向量正交的定义也可推广到内积空间中.

定义 3.7 若 \(f(x), g(x) \in C[a, b]\) 满足 \[
(f, g) = \int_{a}^{b} \rho(x) f(x) g(x) \mathrm{d}x = 0, \tag{3.3.7}
\] 则称 \(f\) 与 \(g\) 在 \([a, b]\) 上带权 \(\rho(x)\) 正交. 若函数族 \(\varphi_{0}(x), \varphi_{1}(x), \cdots, \varphi_{n}(x), \cdots\) 满足关系 \[
(\varphi_{j}, \varphi_{k}) = \int_{a}^{b} \rho(x) \varphi_{j}(x) \varphi_{k}(x) \mathrm{d}x = \begin{cases} 0, & j \neq k, \\ A_{k} > 0, & j = k, \end{cases} \tag{3.3.8}
\]

\href{../../source-images/pdf-067.jpeg}{查看原页}

则称 \(\{\varphi_k\}\) 是 \([a,b]\) 上带权 \(\rho(x)\) 的正交函数族; 若 \(A_k \equiv 1\), 则称 \(\{\varphi_k\}\) 为标准正交函数族. 例如, 三角函数族 \[
1, \cos x, \sin x, \cos 2x, \sin 2x, \cdots
\] 就是在区间 \([-\pi, \pi]\) 上的正交函数族 (权 \(\rho(x) \equiv 1\)), 其 \[
(1,1) = 2\pi,
\] \[
(\sin nx, \sin mx) = \int_{-\pi}^{\pi} \sin nx \sin mx \mathrm{d}x = \begin{cases} \pi, & m=n, \\ 0, & m \neq n \end{cases} (n,m=1,2, \cdots),
\] \[
(\cos nx, \cos mx) = \int_{-\pi}^{\pi} \cos nx \cos mx \mathrm{d}x = \begin{cases} \pi, & m=n, \\ 0, & m \neq n \end{cases} (n,m=1,2, \cdots),
\] \[
(\cos nx, \sin mx) = \int_{-\pi}^{\pi} \cos nx \sin mx \mathrm{d}x = 0 (n,m=1,2, \cdots).
\] 在 \(\mathbf{R}^n\) 空间中, 任一向量都可用它的一组线性无关的基表示, 内积空间的任一元素 \(f(x) \in C[a,b]\) 也同样可用线性无关的基表示, 此时相应地有如下定义.

定义 3.8 设 \(\varphi_0(x), \varphi_1(x), \cdots, \varphi_{n-1}(x)\) 在 \([a,b]\) 上连续, 如果 \[
a_0 \varphi_0(x) + a_1 \varphi_1(x) + \cdots + a_{n-1} \varphi_{n-1}(x) = 0
\] 当且仅当 \(a_0 = a_1 = \cdots = a_{n-1} = 0\) 时成立, 则称 \(\varphi_0, \varphi_1, \cdots, \varphi_{n-1}\) 在 \([a,b]\) 上是线性无关的.

若函数族 \(\{\varphi_k\} (k=0,1, \cdots)\) 中的任何有限个 \(\varphi_k\) 线性无关, 则称 \(\{\varphi_k\}\) 为线性无关函数族.

例如, \(1,\allowbreak  x,\allowbreak  x^2,\allowbreak  \cdots,\allowbreak  x^n,\allowbreak  \cdots\) 就是 \([a,b]\) 上的线性无关函数族. 若 \(\varphi_0(x), \varphi_1(x), \cdots, \varphi_{n-1}(x)\) 是 \([a,b]\) 上的线性无关函数, 且 \(a_0, a_1, \cdots, a_{n-1}\) 是任意实数, 则 \[
S(x) = a_0 \varphi_0(x) + a_1 \varphi_1(x) + \cdots + a_{n-1} \varphi_{n-1}(x)
\] 的全体是 \(C[a,b]\) 中的一个子集, 记作 \[
\Phi = \operatorname{span}\{\varphi_0, \varphi_1, \cdots, \varphi_{n-1}\}. \tag{3.3.9}
\] 下面给出判断函数族 \(\{\varphi_k\} (k=0,1, \cdots, n-1)\) 线性无关的充要条件.

定理 3.6 \(\varphi_0(x), \varphi_1(x), \cdots, \varphi_{n-1}(x)\) 在 \([a,b]\) 上线性无关的充分必要条件是它的 Cramer 行列式 \(G_{n-1} \neq 0\), 其中 \[
G_{n-1} = G(\varphi_0, \varphi_1, \cdots, \varphi_{n-1}) = \left| \begin{array}{cccc} (\varphi_0, \varphi_0) & (\varphi_0, \varphi_1) & \cdots & (\varphi_0, \varphi_{n-1}) \\ (\varphi_1, \varphi_0) & (\varphi_1, \varphi_1) & \cdots & (\varphi_1, \varphi_{n-1}) \\ \vdots & \vdots & & \vdots \\ (\varphi_{n-1}, \varphi_0) & (\varphi_{n-1}, \varphi_1) & \cdots & (\varphi_{n-1}, \varphi_{n-1}) \end{array} \right|. \tag{3.3.10}
\] 定理的证明由读者自行完成.

\begin{center}\rule{0.5\linewidth}{0.5pt}\end{center}

\subsection{相关原页校注（agent 补充）}\label{ux76f8ux5173ux539fux9875ux6821ux6ce8agent-ux8865ux5145}

以下校注来自本节覆盖的原页，可能也涉及同页相邻小节；原文照录与校正说明分开保留。

\subsubsection{原书 PDF 页 67 的校注}\label{ux539fux4e66-pdf-ux9875-67-ux7684ux6821ux6ce8}

\href{../pages/pdf-067.md}{完整原页转录} · \href{../../source-images/pdf-067.jpeg}{原页图像}

校注记录：ch03a-E05。

\begin{quote}
校注（agent 补充，原书术语疑误）：原页定理 3.6 写作``Cramer 行列式''，按原文保留。式 (3.3.10) 是由内积组成的 Gram 行列式，通常称``Gram（格拉姆）行列式''。
\end{quote}

\subsection{本节覆盖原页的校注记录}\label{ux672cux8282ux8986ux76d6ux539fux9875ux7684ux6821ux6ce8ux8bb0ux5f55}

下列记录按原页关联，含同页相邻小节的校注；计算时同时读取上文校注和完整原页。

\begin{itemize}
\tightlist
\item
  \texttt{ch03a-E05}：原书 PDF 页 67
\end{itemize}

\end{document}
